
% 附录源 PDF 第 26—28 页；A.4.4、A.4.5 全文。
\subsection{基于有效影响曲线边际期望的在线估计一阶展开}

现在为在线估计建立相应的展开。

\begin{bookresult}{定理}{A.5}
回顾 $D^N_{g^*}(P)=N^{-1}\sum_{i=1}^N\bar D_{g^*,i}(P)$，其中
\[
\bar D_{g^*,i}(P)(O^N)
 =\sum_{x\in\{y,a,w\}}\sum_{s=1}^{\tau}\bar D_{x(s)}(P)(X(i),C_x(i)),
\]
$\bar D_{x(s)}$ 是有效影响曲线中对应于给定 $C_x(s)$ 时 $X(s)$ 条件分布得分的分量。
有
\[
-P_0^N\bar D_{g^*,i}(P)
 =\Psi_{g^*}(P)-\Psi_{g^*}(P_0)+R_{21,g^*,i}(P,P_0),
\]
其中 $R_{21,g^*,i}(P,P_0)=R_{21a,g^*}(\theta,\theta_0)+R_{21b,g^*,i}(P,P_0)$，且
\[
\begin{aligned}
R_{21a,g^*}(\theta,\theta_0)
={}&\sum_{x\in\{w,a,y\}}\sum_{s=1}^{\tau}
 \left\{\prod_{t=1}^{s-1}P^*_{o(t)}\right\}
 \Delta_x(P^*_{o(s)},P^*_{0,o(s)})\\
&\quad\times\Biggl\{\sum_{j=s+1}^{\tau}
 \prod_{l=s+1}^{j-1}P^*_{o(l)}(P^*_{o(j)}-P^*_{0,o(j)})\\
&\hspace{12em}\times\prod_{l=j+1}^{\tau}P^*_{0,o(l)}\Biggr\}Y(\tau),
\end{aligned}
\]
\[
R_{21b,g^*,i}(P,P_0)
 =\sum_{x\in\{w,a,y\}}\sum_{s=1}^{\tau}R_{21b,x(s),g^*,i}(P,P_0),
\]
\[
\begin{aligned}
&R_{21b,x(s),g^*,i}(P,P_0)\\
&=\int_c\int_x\left(\frac{h_{c_x(i),0}}{\bar h_{c_x}}(c)-1\right)
 \Phi_{x(s)}(x,c)\\
&\qquad\times\mathrm d(P^*_{x(s)}-P^*_{0,x(s)})(x\mid c)
 h^*_{c_x(s)}(c)\,\mathrm d\mu(c).
\end{aligned}
\]
具体而言，这给出如下恒等式：
\[
\begin{aligned}
-\frac1N\sum_{i=1}^NP_0^N\bar D_{g^*,i}(\theta_{i-1})
={}&\frac1N\sum_{i=1}^N\Psi_{g^*}(\theta_{i-1})-\Psi_{g^*}(\theta_0)\\
&+\bar R_{21,g^*,N},
\end{aligned}
\]
其中
\begin{equation}
\bar R_{21,g^*,N}\equiv\bar R_{21a,g^*,N}+\bar R_{21b,g^*,N},
\tag{A.14}
\end{equation}
且
\[
\begin{aligned}
\bar R_{21a,g^*,N}&\equiv\frac1N\sum_{i=1}^NR_{21a,g^*}(\theta_{i-1},\theta_0),\\
\bar R_{21b,g^*,N}&\equiv\frac1N\sum_{i=1}^NR_{21b,g^*,i}(\theta_{i-1},\theta_0).
\end{aligned}
\]
\end{bookresult}

\begin{proof}
在定理 A.3 的证明中，已得到
\[
\begin{aligned}
&\Psi_{g^*}(\theta)-\Psi_{g^*}(\theta_0)\\
&=\sum_{x\in\{w,a,y\}}\sum_{s=1}^{\tau}
 \bigl(\Psi_{g^*,P,x(s)}(P^*_{x(s)})
 -\Psi_{g^*,P,x(s)}(P^*_{0,x(s)})\bigr)
 +R_{21a,g^*}(\theta,\theta_0).
\end{aligned}
\]
我们将证明
\begin{equation}
\begin{aligned}
-P_0^N\bar D_{x(s),i}(P)
={}&\Psi_{g^*,P,x(s)}(P^*_{x(s)})
 -\Psi_{g^*,P,x(s)}(P^*_{0,x(s)})\\
&+R_{21b,x(s),g^*,i}(P,P_0).
\end{aligned}
\tag{A.15}
\end{equation}
由此得到以下推导：
\[
\begin{aligned}
&-P_0^N\bar D_{g^*,i}(P)\\
&=\sum_{x\in\{y,a,w\}}\sum_{s=1}^{\tau}-P_0^N\bar D_{x(s),i}(P)\\
&=\sum_{x\in\{y,a,w\}}\sum_{s=1}^{\tau}
 \Bigl\{\Psi_{g^*,P,x(s)}(P^*_{x(s)})
 -\Psi_{g^*,P,x(s)}(P^*_{0,x(s)})
 +R_{21b,x(s),g^*,i}(P,P_0)\Bigr\}\\
&=\sum_{x\in\{y,a,w\}}\sum_{s=1}^{\tau}
 \bigl\{\Psi_{g^*,P,x(s)}(P^*_{x(s)})
 -\Psi_{g^*,P,x(s)}(P^*_{0,x(s)})\bigr\}\\
&\quad+\sum_{x\in\{y,a,w\}}\sum_{s=1}^{\tau}R_{21b,x(s),g^*,i}(P,P_0)\\
&=\Psi_{g^*}(P)-\Psi_{g^*}(P_0)-R_{21a,g^*}(\theta,\theta_0)\\
&\quad+\sum_{x\in\{y,a,w\}}\sum_{s=1}^{\tau}R_{21b,x(s),g^*,i}(P,P_0)\\
&\equiv\Psi_{g^*}(P)-\Psi_{g^*}(P_0)-R_{21a,g^*}(\theta,\theta_0)
 +R_{21b,g^*,i}(P,P_0),
\end{aligned}
\]
从而证明了定理。因此，只需再建立式 (A.15)。有
\[
\begin{aligned}
&-P_0^N\bar D_{x(s),i}(P)\\
&=-P_0^N\frac{h^*_{c_x(s)}(C_x(i))}{\bar h_{c_x}(C_x(i))}\\
&\quad\times\Bigl\{E(Y_{g^*}(\tau)\mid X(s)=X(i),C_x(s)=C_x(i))\\
&\hspace{7em}-E(Y_{g^*}(\tau)\mid C_x(s)=C_x(i))\Bigr\}\\
&=P_0^N\frac{h^*_{c_x(s)}(C_x(i))}{\bar h_{c_x}(C_x(i))}
 \int_x\Phi_{x(s)}(x,C_x(i))\\
&\hspace{8em}\times\mathrm d(P^*_{x(s)}-P^*_{0,x(s)})(x\mid C_x(i))\\
&=\int_c\frac{h^*_{c_x(s)}(c)}{\bar h_{c_x}(c)}
 \int_x\Phi_{x(s)}(x,c)\\
&\hspace{6em}\times\mathrm d(P^*_{x(s)}-P^*_{0,x(s)})(x\mid c)
 h_{c_x(i),0}(c)\,\mathrm d\mu(c)\\
&=\int_c\int_x\frac{h_{c_x(i),0}(c)}{\bar h_{c_x}(c)}
 \Phi_{x(s)}(x,c)\,\mathrm d(P^*_{x(s)}-P^*_{0,x(s)})(x\mid c)
 h^*_{c_x(s)}(c)\,\mathrm d\mu(c)\\
&=\int_c\int_x\left(\frac{h_{c_x(i),0}}{\bar h_{c_x}}(c)-1\right)
 \Phi_{x(s)}(x,c)\,\mathrm d(P^*_{x(s)}-P^*_{0,x(s)})(x\mid c)
 h^*_{c_x(s)}(c)\,\mathrm d\mu(c)\\
&\quad+\int_{x,c}\Phi_{x(s)}(x,c)
 \,\mathrm d(P^*_{x(s)}-P^*_{0,x(s)})(x\mid c)
 h^*_{c_x(s)}(c)\,\mathrm d\mu(c)\\
&\equiv R_{21b,x(s),g^*,i}(P,P_0)
 +\Psi_{g^*,P,x(s)}(P^*_{x(s)})
 -\Psi_{g^*,P,x(s)}(P^*_{0,x(s)}).
\end{aligned}
\]
证明完毕。
\end{proof}

\paragraph{余项 $\bar R_{21b,g^*,N}$ 的讨论：另一个平稳性假设。}
尽管外观相似，来自在线一步估计量和在线 TMLE 分析的余项
$\bar R_{21b,g^*,N}$，与来自常规 TMLE 和一步估计量分析的余项
$R_{21b,g^*}$，实际上差别很大。仔细考察可以发现，
$\bar R_{21b,g^*,N}$ 本质上是对 $1\le i\le N$ 的各个 $i$ 特异乘积求平均，乘积形如
\[
\left(1-\frac{h_i(\theta_0)}{N^{-1}\sum_{j=1}^Nh_j(\theta_i)}\right)
 \times(f(\theta_i)-f(\theta_0)),
\]
而 $R_{21b,g^*}$ 对 $i$ 所平均的是
\[
\left(1-\frac{N^{-1}\sum_{j=1}^Nh_j(\theta_0)}
 {N^{-1}\sum_{j=1}^Nh_j(\theta_N)}\right)
 \times(f(\theta_N)-f(\theta_0)).
\]
这一描述虽然粗略，却能说明问题：$f$ 和 $h_i$ 都把 $\theta$ 连续地映射为函数，且 $\theta_N$ 收敛到 $\theta_0$。
看起来，$\bar R_{21b,g^*,N}$ 只从 $f(\theta_N)-f(\theta_0)$ 趋于零中获益；
而 $R_{21b,g^*}$ 则通过乘积，同时利用了 $f(\theta_N)-f(\theta_0)$ 和
\[
1-\frac{N^{-1}\sum_{j=1}^Nh_j(\theta_0)}{N^{-1}\sum_{j=1}^Nh_j(\theta_N)}
\]
这两个因子趋于零的性质。这说明，$\bar R_{21b,g^*,N}$ 的行为可能明显差于 $R_{21b,g^*}$。

还需要对 $\bar R_{21b,g^*,N}$ 作更深入的研究。这里指出，在时间序列满足较强平稳性时，可以把 $\bar R_{21b,g^*,N}$ 表达式中与 $i$ 对应的项换为
\[
\left(1-\frac{N^{-1}\sum_{j=1}^Nh_j(\theta_0)}
 {N^{-1}\sum_{j=1}^Nh_j(\theta_i)}\right)
 \times(f(\theta_i)-f(\theta_0)),
\]
于是，当 $i,N$ 趋于无穷时，$\bar R_{21b,g^*,N}$ 也可以同时利用两个因子趋于零的性质。
具体而言，仅假设给定 $\bar O(t-1)$ 时产生 $O(t)$ 的机制不随时间 $t$ 改变，可能还不够；可以再增加一个假设：对于 $x=w,a,y$，$C_x(i)$ 的边际密度不随 $i$ 改变。

\subsection{基于有效影响曲线条件期望的在线估计一阶展开}

下面的定理可由定理 A.5 直接推出。

\begin{bookresult}{定理}{A.6}
有
\[
\begin{aligned}
-\frac1N\sum_{i=1}^NP_{0,o(i)}\bar D_{g^*}(\theta_{i-1})
={}&\frac1N\sum_{i=1}^N\Psi_{g^*}(\theta_{i-1})-\Psi_{g^*}(\theta_0)\\
&+\bar R_{21,g^*,N}+\bar R_{22,g^*,N},
\end{aligned}
\]
其中 $\bar R_{21,g^*,N}$ 由式 (A.14) 定义，且
\[
\bar R_{22,g^*,N}\equiv\frac1N\sum_{i=1}^N(P_{0,o(i)}-P_0^N)
 \{\bar D_{g^*}(\theta_{i-1})-\bar D_{g^*}(\theta_0)\}.
\]
\end{bookresult}
